\documentclass{article}
\usepackage{hyperref}
\hypersetup{
    colorlinks=true,
    linkcolor=blue,
    filecolor=magenta,
    urlcolor=cyan
}
\title{HTML Notes}
\author{Zachary Sears}
\begin{document}
\maketitle
\section{General Notes} \label{sec:General}
From the Javascript tutorial:\\
"Where do I Put my Javascript? How to link Javascript
to HTML \textbar\ Tutorial for Beginners":
The script element: <script></script>, similarly to the link element, allows us to:
\begin{itemize}
    \item Write JavaScript directly into the element
    \item Link a src JavaScript file
\end{itemize}

\section{Lesson 1} \label{sec:L1}
Playlist for learning HTML is found \href{https://www.youtube.com/playlist?list=PL0Zuz27SZ-6OlAwitnFUubtE93DO-l0vu}{here}.
All HTML files are a combination of element tags and text.

Each file begins with an \verb|<html>| opening tag and ends with an \verb|</html>| closing tag.
Within these two tags is the actual webpage definition. The first half of which is the \verb|<head></head>| tag set which holds
and defines data not on the actual page. However, if a \verb|<title></title>| element is defined, this will show up in the page's
tab at the top of the browser. The second half of the page is the \verb|<body></body>| tag set which defines what the page displays.
There are heading tags, such as \verb|<h1></h1>| which displays header size text, and paragraph tags: \verb|<p></p>| which dispalys
blocks of text.
Webpage validation can be done at \url{https://validator.w3.org}. This site allows file uploads to check against current HTML conventions
and definitions.
\section{Lesson 2} \label{sec:L2}
An author and page description tag in the head of an HTML index file is useful for seach engines to learn more about a page. These are not
displayed on the page itself and are formatted as follows:
\begin{itemize}
    \item \verb|<meta name = "author" content = "...">|
    \item \verb|<meta name = "description" content = "...">|
\end{itemize}
There is a \verb|<link>| element that can be placed in the \verb|<head>| element. \verb|<link>| has three attributes:
\begin{enumerate}
    \item rel (relation) - How this element relates to the page.
    \item href (hypertext reference) - A reference to a resource that will be used by the page.
    \item type - What kind of resource is being referenced by this link.
\end{enumerate}
CSS code can be relegated to a separate file and replace with a link to that file using the previously described link element.
\section{Lesson 3} \label{sec:L3}
Moved the the CSS and image files into the parent folder and the links broke.
Looked up a fix for this and found that because
the files are in the parent folder, I can use the \verb|../<filename>| relative path to reference them.\\
\\
New tags:
\begin{itemize}
    \item \verb|<hr>| - Horizontal rule, places a horizontal line across the page.
    \item \verb|<br>| - Line break, inserts a line break.
    \item \verb|<em>| - Emphasis, emphasizes the text inside the element. Does not create a new line.
    \item \verb|<strong>| - Strong emphasis, also emphasizes the text but bolds the text instead of itallicizing it.
    \item \verb|&| - Begins an HTML entity. A list of HTML entities can be found \href{https://html.spec.whatwg.org/multipage/named-characters.html#named-character-references}{here}
    \begin{itemize}
        \item \verb|&nbsp| - Adds a space into a text block.
        \item \verb|&lt| - Less than: \verb|>|
        \item \verb|&gt| - Greater than: \verb|<|
        \item \verb|&copy| - Interts the copyright symbol.
    \end{itemize}
    \item \verb|<address>| - Displays an address on the page.
    \item \verb|<!--...-->| - Comment structure.
\end{itemize}

\section{Lesson 4} \label{sec:L4}
More new elements!
\begin{itemize}
    \item \verb|<ol>| - Ordered list, create a numbered list by default.
    \item \verb|<ul>| - Unordered list, create a bulleted list by default.
    \item \verb|<li>| - List item
    \item \verb|<dl>| - Description list, create a list of items with indented descriptions.
    \begin{itemize}
        \item \verb|<dt>| - Description term, name of the list item.
        \item \verb|<dd>| - Description details - details about the description term.
        \item Format of \verb|<dl>| in a webpage by default is:
        \begin{tabbing}
            \verb|<body>| \\ 
            \hspace{1em} \= \verb|<dl>| \\
            \hspace{1em} \= \hspace{1em} \= item 1 \\
            \hspace{1em} \= \hspace{1em} \= \hspace{1em} \= item 1 details \\
            \hspace{1em} \= \hspace{1em} \= item 2 \\
            \hspace{1em} \= \hspace{1em} \= \hspace{1em} \= item 2 details \\
            \hspace{1em} \= \verb|</dl>| \\
            \verb|</body>|\\
        \end{tabbing}
    \end{itemize}
\end{itemize}

\section{Lesson 5} \label{sec:L5}
\verb|<a>| - Anchor tage, used to link other webpages. Attributes for href are defined in \ref{sec:L2} under the
\verb|<link>| element.
\end{document}